%%%% Настройки автора 
%% 
%% 	 Пожалуйста, ознакомьтесь с функционалом шаблона из [1,2], а также с пакетами, подключенными в ch_preamble.
%% 
%%   Новым командам лучше присваивать уникальные имена.
%% 
%%%% Author settings
%% 
%%   Please, see all possible packages using the search in files of ch_preamble. 
%%   
%%   Please, for user-defined commands write only unique command titles.
%%


%%%% Настройки подсветки синтаксиса листингов
\lstset{
  keywordstyle=\color{blue}\bfseries,
  emphstyle=\color{Mulberry},
}

%%%% Определение языка HCL для листингов
\lstdefinelanguage{HCL}{
  morecomment=[l]{\#},
  morestring=[b]",
  morekeywords={resource, variable, output, module, provider,
    for_each, locals, data, true, false, null},
}

%%%% Дополнительные ключевые слова для Go
\lstdefinelanguage{Go}{
  morekeywords={break, default, func, interface, select,
    case, defer, go, map, struct, chan, else, goto, package,
    switch, const, fallthrough, if, range, type, continue,
    for, import, return, var, nil, true, false, string, int,
    bool, error, byte, make, len, append, new, context},
  morecomment=[l]{//},
  morecomment=[s]{/*}{*/},
  morestring=[b]",
  morestring=[b]`,
}

%%%% Нормоконтроль: текст внутри таблиц набирается 12 pt (основной текст — 14 pt)
%% \AtBeginEnvironment из etoolbox (подключён в шаблоне). Подпись «Таблица N — …»
%% у float-таблиц (\caption вне tabular) остаётся в основном кегле.
\AtBeginEnvironment{tabular}{\fontsize{12pt}{14pt}\selectfont}
\AtBeginEnvironment{tabularx}{\fontsize{12pt}{14pt}\selectfont}
\AtBeginEnvironment{longtable}{\fontsize{12pt}{14pt}\selectfont}

%%%% Подпись к листингу-рисунку (код оформлен как рисунок, см. главы 2–4).
%% \captionof{figure} вне float применяет настройки [figure] (обнуление
%% \parindent, labelformat=figwithoutspace и пр.) прямо в основном потоке, и
%% они «утекают» в последующий текст: пропадает красная строка, а подписи
%% таблиц теряют пробел («Таблица3.1»). Оборачиваем подпись в группу {…},
%% чтобы все локальные настройки caption остались внутри неё.
%% Счётчик figure и \label при этом работают штатно (refstepcounter глобален).
\makeatletter
\newcommand{\codefig}[2]{%
  \par
  \begingroup
    \captionof{figure}{#1}\label{#2}%
  \endgroup
  \par
}
\makeatother

%%%% Листинг-рисунок как единый неразрывный блок (minipage): код + подпись не
%% разрываются границей страницы и в них не вклинивается плавающая таблица/рисунок.
%% Если блок не помещается в остаток страницы — целиком переносится на следующую.
%% (Все листинги основного текста короче высоты полосы, поэтому переполнения нет.)
\newenvironment{codeblock}%
  {\par\addvspace{\mfloatsep}\noindent\minipage{\linewidth}}%
  {\endminipage\par\addvspace{\mfloatsep}}


%%%% Подключение библиографии / Upload bibliography
%%
%%
\addbibresource{chapters/013_references.bib} %



%%%% Полезные настройки / Usefull settings
%% 
%% Раскомментируйте, чтобы
%%
%% pdf при открытии выравнивался по окну
%% pdf fit screen window
\hypersetup{
pdfstartview={FitBH}
}
%% перенумеровать все строки pdf
%% enumerate all lines in pdf 
%\usepackage{lineno}
%\linenumbers
%%
%% установить дату после названия ВКР - расскоментируйте код в title.tex
%% set data after the thesis title - uncomment code in title.tex
\let\ordinal\relax %avoid extra warning
\usepackage{datetime}



%% In case of deleting the following info, please, delete the examples in the chapter body.

%% В случае комментирования (удаления) следующего кода могут появиться ошибки при компиляции примеров, т.е. необходимо будет удалить и примеры в теле главы.

\newcommand{\overbar}[1]{\mkern 1.5mu\overline{\mkern-1.5mu#1\mkern-1.5mu}\mkern 1.5mu}

%http://tex.stackexchange.com/questions/16645/blackboard-italic-font
% for itallic sign of context K to be a parametr
\DeclareMathAlphabet{\mathbbmsl}{U}{bbm}{m}{sl}
\newcommand{\cont}[1][K]{\ensuremath{\mathbbmsl{#1}}}

%%ARROWS

%mu = math unit = 1em
%\mkern-18mu
%"minus quad"

%https://tex.stackexchange.com/a/389805/44348
\newcommand{\fcaarrow}[1]{%
	{}^{\scriptscriptstyle\bm{#1}}
}
%%%%%%%%%%%%%%%%%%%%%%% ARROWS from Formal Concept Analysis
% small and bold \uparrow
\newcommand{\uA}{\fcaarrow{{\uparrow\mkern-12mu}}}
% small and bold \downarrow
\newcommand{\dA}{\fcaarrow{\downarrow\mkern-2mu}}
% small and bold \uparrow+\downarrow
\newcommand{\ud}{\fcaarrow{\uparrow\mkern-12mu}\fcaarrow{\downarrow\mkern-2mu}}
% small and bold \downarrow+\uparrow
\newcommand{\du}{\fcaarrow{\downarrow\mkern-2mu}\fcaarrow{\uparrow\mkern-12mu}}


%http://tex.stackexchange.com/questions/74125/how-do-i-put-text-over-symbols
\newcommand\eqdef{\mathrel{\overset{\makebox[0pt]{\mbox{\normalfont\tiny def}}}{=}}} %\sffamily



%%% Правила задания нового окружения

\theoremstyle{myplain} % первая команда для ввода доказательств
\newtheorem{m-new-env-first}{Название\_окружения}[chapter] 
% вместо m-new-env-first необходимо подставить название нового окружения;
% вместо Название\_окружения необходимо подставить название окружения, выводящееся в pdf;
% последний параметр обеспечивает нумерацию в пределах главы не меняется


\theoremstyle{mydefinition} % первая команда для ввода окружений, не связанных с доказательствами
\newtheorem{m-new-env-second}{Название\_окружения}[chapter] 
% вместо m-new-env-second необходимо подставить название нового окружения;
% вместо Название\_окружения необходимо подставить название окружения, выводящееся в pdf;
% последний параметр обеспечивает нумерацию в пределах главы не меняется